\subsection{Adaptive Signal Smoothing: The \texorpdfstring{1\euro{}}{1-Euro} Filter}
\label{subsec:one_euro_filter}

Raw tracking signals are inherently noisy, producing \emph{jitter} ---
rapid, small-amplitude oscillations around the true value that make
a stationary target appear to vibrate.
The standard remedy, a low-pass filter, suppresses high-frequency
noise but introduces \emph{lag}: the filtered output responds
sluggishly to genuine motion, reducing interactivity.
Casiez et al.\ \cite{casiez2012oneeuro} frame this as a fundamental
trade-off: empirical studies show that lag above 60\,ms and jitter
above 1\,mm mean-to-peak both measurably degrade user performance,
leaving only a narrow acceptable window for any filter to operate in.

The 1\euro{} filter resolves the trade-off by making the cutoff
frequency $f_c$ adaptive.
Rather than applying the same degree of smoothing at all times, it
scales $f_c$ linearly with the magnitude of the signal's instantaneous
speed \cite{casiez2012oneeuro}:
%
\begin{equation}
    f_c \;=\; f_{c_{\min}} \;+\; \beta\,\bigl|\dot{\hat{X}}_i\bigr|
    \label{eq:one_euro}
\end{equation}
%
At low speeds $f_c$ approaches $f_{c_{\min}}$, aggressively
attenuating jitter.
As speed increases, $f_c$ rises proportionally, keeping the response
tight and minimising lag.
This matches the asymmetry of human perception: people are most
sensitive to jitter during slow or stationary motion, and most
sensitive to lag during fast motion \cite{casiez2012oneeuro}.

The filter exposes exactly two parameters.
$f_{c_{\min}}$ controls slow-speed smoothing: lowering it reduces
jitter at rest.
$\beta$ controls high-speed responsiveness: raising it prevents the
filter from lagging behind fast motion.
The two are independent and can be tuned sequentially without any
background in signal processing.
In a direct comparison against moving-average, exponential-smoothing,
and Kalman-filter alternatives, the 1\euro{} filter achieves the
lowest overall error and the least lag across all speed intervals
\cite{casiez2012oneeuro}, making it well suited to real-time
landmark smoothing where both artefacts must be minimised
simultaneously.
